%%% Hlavní soubor. Zde se definují základní parametry a odkazuje se na ostatní části. %%%

%% Verze pro jednostranný tisk:
% Okraje: levý 40mm, pravý 25mm, horní a dolní 25mm
% (ale pozor, LaTeX si sám přidává 1in)
\documentclass[12pt,a4paper]{report}
\setlength\textwidth{145mm}
\setlength\textheight{247mm}
\setlength\oddsidemargin{15mm}
\setlength\evensidemargin{15mm}
\setlength\topmargin{0mm}
\setlength\headsep{0mm}
\setlength\headheight{0mm}
% \openright zařídí, aby následující text začínal na pravé straně knihy
\let\openright=\clearpage

%% Pokud tiskneme oboustranně:
% \documentclass[12pt,a4paper,twoside,openright]{report}
% \setlength\textwidth{145mm}
% \setlength\textheight{247mm}
% \setlength\oddsidemargin{15mm}
% \setlength\evensidemargin{0mm}
% \setlength\topmargin{0mm}
% \setlength\headsep{0mm}
% \setlength\headheight{0mm}
% \let\openright=\cleardoublepage

%% Pokud používáte csLaTeX (doporučeno):
\usepackage{czech}
%% Pokud nikoliv:
%\usepackage[czech]{babel}
%\usepackage[T1]{fontenc}

%% Použité kódování znaků: obvykle latin2, cp1250 nebo utf8:
\usepackage[utf8]{inputenc}

%% Ostatní balíčky
\usepackage{graphicx}
\usepackage{amsthm}

%% Balíček hyperref, kterým jdou vyrábět klikací odkazy v PDF,
%% ale hlavně ho používáme k uložení metadat do PDF (včetně obsahu).
%% POZOR, nezapomeňte vyplnit jméno práce a autora.
\usepackage[ps2pdf,unicode]{hyperref}   % Musí být za všemi ostatními balíčky
\hypersetup{pdftitle=Hledání pěších cest v mapě}
\hypersetup{pdfauthor=Tomáš Pokorný}

%%% Drobné úpravy stylu

% Tato makra přesvědčují mírně ošklivým trikem LaTeX, aby hlavičky kapitol
% sázel příčetněji a nevynechával nad nimi spoustu místa. Směle ignorujte.
\makeatletter
\def\@makechapterhead#1{
  {\parindent \z@ \raggedright \normalfont
   \Huge\bfseries \thechapter. #1
   \par\nobreak
   \vskip 20\p@
}}
\def\@makeschapterhead#1{
  {\parindent \z@ \raggedright \normalfont
   \Huge\bfseries #1
   \par\nobreak
   \vskip 20\p@
}}
\makeatother

% Toto makro definuje kapitolu, která není očíslovaná, ale je uvedena v obsahu.
\def\chapwithtoc#1{
\chapter*{#1}
\addcontentsline{toc}{chapter}{#1}
}

\begin{document}

% Trochu volnější nastavení dělení slov, než je default.
\lefthyphenmin=2
\righthyphenmin=2

%%% Titulní strana práce

\pagestyle{empty}
\begin{center}

\large

Univerzita Karlova v Praze

\medskip

Matematicko-fyzikální fakulta

\vfill

{\bf\Large BAKALÁŘSKÁ PRÁCE}

\vfill

\centerline{\mbox{\includegraphics[width=60mm]{../img/logo.eps}}}

\vfill
\vspace{5mm}

{\LARGE Tomáš Pokorný}

\vspace{15mm}

% Název práce přesně podle zadání
{\LARGE\bfseries Hledání pěších cest v mapě}

\vfill

% Název katedry nebo ústavu, kde byla práce oficiálně zadána
% (dle Organizační struktury MFF UK)
Katedra aplikované matematiky

\vfill

\begin{tabular}{rl}

Vedoucí bakalářské práce: & Mgr. Martin Mareš, Ph.D. \\
\noalign{\vspace{2mm}}
Studijní program: & Informatika \\
\noalign{\vspace{2mm}}
Studijní obor: & Obecná informatika \\
\end{tabular}

\vfill

% Zde doplňte rok
Praha 2014

\end{center}

\newpage

%%% Následuje vevázaný list -- kopie podepsaného "Zadání bakalářské práce".
%%% Toto zadání NENÍ součástí elektronické verze práce, nescanovat.

%%% Na tomto místě mohou být napsána případná poděkování (vedoucímu práce,
%%% konzultantovi, tomu, kdo zapůjčil software, literaturu apod.)

\openright

\noindent
Poděkování.
- české komunitě OSM
- vedoucímu
- Gaventovi

\newpage

%%% Strana s čestným prohlášením k bakalářské práci

\vglue 0pt plus 1fill

\noindent
Prohlašuji, že jsem tuto bakalářskou práci vypracoval(a) samostatně a výhradně
s~použitím citovaných pramenů, literatury a dalších odborných zdrojů.

\medskip\noindent
Beru na~vědomí, že se na moji práci vztahují práva a povinnosti vyplývající
ze zákona č. 121/2000 Sb., autorského zákona v~platném znění, zejména skutečnost,
že Univerzita Karlova v Praze má právo na~uzavření licenční smlouvy o~užití této
práce jako školního díla podle §60 odst. 1 autorského zákona.

\vspace{10mm}

\hbox{\hbox to 0.5\hsize{%
V ........ dne ............
\hss}\hbox to 0.5\hsize{%
Podpis autora
\hss}}

\vspace{20mm}
\newpage

%%% Povinná informační strana bakalářské práce

\vbox to 0.5\vsize{
\setlength\parindent{0mm}
\setlength\parskip{5mm}

Název práce:  	
Hledání pěších cest v mapě
% přesně dle zadání

Autor:
Tomáš Pokorný

Katedra:  % Případně Ústav:
Katedra aplikované matematiky
% dle Organizační struktury MFF UK

Vedoucí bakalářské práce:
Mgr. Martin Mareš, Ph.D., Katedra aplikované matematiky
% dle Organizační struktury MFF UK, případně plný název pracoviště mimo MFF UK

Abstrakt:
Cestování pěšky po městě je součástí každodenního života mnoha lidí. Najít
vhodnou cestu nemusí být ve městě jednoduché. Snažili jsme se vytvořit z dat
projektu OpenStreetMap datové struktury vhodné pro vyhledávání pěších tras ve
městě, včetně možnosti chůze mimo cesty v průchozích oblastech (parky, náměstí). 
Součástí práce je i řešení některých druhů chyb ve vstupních datech. Vytvořili
jsme také aplikaci umožňující trasy hledat a exportovat do formátu GPX.


% abstrakt v rozsahu 80-200 slov; nejedná se však o opis zadání bakalářské práce

Klíčová slova:
mapa, pěší vzdálenosti, hledání trasy
% 3 až 5 klíčových slov

\vss}\nobreak\vbox to 0.49\vsize{
\setlength\parindent{0mm}
\setlength\parskip{5mm}

Title:
Finding footpaths in a map
% přesný překlad názvu práce v angličtině

Author:
Tomáš Pokorný

Department:
Department of Applied Mathematics
% dle Organizační struktury MFF UK v angličtině

Supervisor:
Mgr. Martin Mareš, Ph.D., Department of Applied Mathematics
% dle Organizační struktury MFF UK, případně plný název pracoviště
% mimo MFF UK v angličtině

Abstract:
Walking in the city is a part of everyday life for many people. Finding the right
way is not very easy. From the data of OpenStreetMap project we made data
structures suitable for finding paths for pedesterians, including direct walking through
areas (parks, squares). Fixing some bugs in input data is also a part of the
thesis. We have made an application for finding pedesterian ways and their export
to GPX as well.

% abstrakt v rozsahu 80-200 slov v angličtině; nejedná se však o překlad
% zadání bakalářské práce

Keywords:
map, walking distance, finding paths
% 3 až 5 klíčových slov v angličtině

\vss}

\newpage

%%% Strana s automaticky generovaným obsahem bakalářské práce. U matematických
%%% prací je přípustné, aby seznam tabulek a zkratek, existují-li, byl umístěn
%%% na začátku práce, místo na jejím konci.

\openright
\pagestyle{plain}
\setcounter{page}{1}
\tableofcontents

%%% Jednotlivé kapitoly práce jsou pro přehlednost uloženy v samostatných souborech
\chapwithtoc{Úvod}

chceme hledat cestu
chceme chodit pěšky
existují různé služby
většina určena primárně pro auta -> nereflektuje pěší potřeby
 - kros průchozích prostranství


Hledání cesty ve městě je častou situací většiny lidí. Ve městě se vyskytuje
mnoho různých překážek jako jsou ploty, zábradlí či frekventované silnice,
i mnoho průchozích prostranství, například náměstí, parky či sady. V naší práci
se snažíme z mapových dat vytvořit formát vhodný pro rychlé vyhledávání pěších
tras využívajících i průchozí prostranství. Tato práce by měla být jedním z~modulů
budoucí aplikace pro vyhledávání spojení pěšky a MHD po městě. 

%\section{Co chci počítat}
Abychom mohli rychle vyhledávat pěší trasy, potřebujeme k~tomu mít mapu vhodně
reprezentovanou. Vstupní mapová data postupně zpracováváme a vybíráme z nich
použitelné informace. Na konci tohoto procesu vytvoříme graf popisující možné
pěší trasy. V tomto grafu již můžeme vyhledávat pomocí grafových algoritmů pro
hledání cest, v ukázkové aplikaci využíváme Dijkstrův algoritmus. 

%\section{Zdroje dat}
Abychom mohli vyhledávat trasy, potřebujeme k~tomu vhodné mapové podklady.
Protože jsme chtěli, aby bylo možné zpracovaná data volně používat a šířit, 
potřebovali jsme získat i takové mapové podklady. Proto jsme si vybrali jako
zdroj dat OpenStreetMap (OSM), volně dostupné mapové podklady vytvářené komunitou. 
Projekt OpenStreetMap využívá k~tvorbě mapy mimo práce dobrovolníků také jiné
volně dostupné mapové podklady a poskytuje pravěpodobně nejlepší veřejně
dostupná mapová data. Jako zdroj dat o~nadmořské výšce jsme použili data
z~projektu SRTM, která jsou také volně dostupná.

%\section{Co už kdo napsal}
K~porovnání výsledků naší aplikace jsme použili nejznámější webové mapové
aplikace -- Google Maps a Mapy.cz. Také jsme nalezené trasy porovnávali s~jinými
vyhledávači používajícími data OSM -- OsmAnd a TODO. Nalezené trasy jsme také
porovnávali s~vlastní znalostí terénu a skutečně používaných cest.

%TODO: obrazek - vstupni data a vystupni data

V první kapitole se zabýváme zdrojovými geografickými daty. Popíšeme, jakým
způsobem OpenStreetMap popisuje objekty reálného světa a jaké to přináší
problémy při zpracování. Ukážeme jeden ze způsobů, kterým jsou data ukládána -
OSM XML. Popíšeme, které informace z dat OSM využíváme a s jakými problémy jsme
se při jejich zpracování museli vypořádat. Na konci se zmíníme o používaných
výškových datech z projektu SRTM.

Ve druhé kapitole popíšeme používané formáty. Popíšeme souřadný systém UTM a
důvody, proč jsme ho zvolili jako interní reprezentaci souřadnic. Popíšeme
ProtocolBuffery používané jako binární reprezentace našich datových struktur.
Popíšeme datovou strukturu používanou při zpracovávání mapových dat a strukturu
vyhledávacího grafu z dat vzniklého.

Ve třetí kapitole popíšeme použité algoritmy

Ve čtvrté kapitole poíšeme detaily implementace jednotlivých fází zpracování,
popíšeme problémy, které při implmentaci nastaly a jakým způsobem jsme je
řešili. 

V páté kapitule se budeme zabývat hledáním ve vygenerovaných datech. Popíšeme
Dijkstrův algoritmus pro vyhledávání v grafu a popíšeme ukázkovou vyhledávací
aplikaci. 

V šesté kapitole porovnáme výsledky naší aplikace s výsledky jiných vyhledávačů
a vlastními zkušenostmi. 



\include{kap1}
\include{kap2}

% Ukázka použití některých konstrukcí LateXu (odkomentujte, chcete-li)
% \include{ukazka}

\chapter*{Závěr}
\addcontentsline{toc}{chapter}{Závěr}

zrychleni - potencialova heuristika, viz MJ - skripta z GA

skriptovani vypoctu penalt

Další možné rozšíření systému penalt by mohlo být například posuzování intervalů
linek, kdy bychom mohli preferovat častěji jedoucí linky. 

posun odchodu podle prvniho MHD

okolí cest


%%% Seznam použité literatury
%%% Seznam použité literatury je zpracován podle platných standardů. Povinnou citační
%%% normou pro bakalářskou práci je ISO 690. Jména časopisů lze uvádět zkráceně, ale jen
%%% v kodifikované podobě. Všechny použité zdroje a prameny musí být řádně citovány.

\def\bibname{Seznam použité literatury}
\begin{thebibliography}{99}
\addcontentsline{toc}{chapter}{\bibname}

\bibitem{lamport94}
  {\sc Lamport,} Leslie.
  \emph{\LaTeX: A~Document Preparation System}.
  2. vydání.
  Massachusetts: Addison Wesley, 1994.
  ISBN 0-201-52983-1.

\bibitem{osmweb}
	\emph{OpenStreetMap} [online]. 
	2014, 4.5.2014 [cit. 2014-05-04]. 
	Dostupné z: \url{http://www.openstreetmap.org}

\bibitem{osmxml}
	\emph{OSM XML} [online]. 
	2014, 16 September 2013 [cit. 2014-05-06]. 
	Dostupné z: \url{http://wiki.openstreetmap.org/wiki/OSM_XML}

\bibitem{osmfeatures}
	\emph{Cs: Map Features} [online]. 
	2014-04-01 [cit. 2014-05-17]. 
	Dostupné z: \url{http://wiki.openstreetmap.org/wiki/Cs:Map_Features}

\bibitem{wgsnorma}
	{\sc NIMA.}
	\emph{World geodetic system 1984} [online]. 
	3. vydání, 1. dodatek.
	3 January 2000 [cit. 2014-05-04]. 
	Dostupné z: \url{http://www2.jpl.nasa.gov/srtm/}.

\bibitem{utmnorma}
	{\sc Snyder,} John P.
	\emph{Map Projections: A~Working Manual} [online]. 
	1987, [cit. 2014-05-08]. 
	Dostupné z: \url{http://pubs.er.usgs.gov/publication/pp1395}.

\bibitem{srtmweb}
	{\sc NASA.}
	\emph{The Shuttle Radar Topography Mission} [online]. 
	2014, June 17, 2009 [cit. 2014-05-04]. 
	Dostupné z: \url{http://www2.jpl.nasa.gov/srtm/}.

\bibitem{pbfweb}
	{\sc Google.}
	\emph{Protocol Buffers} [online]. 
	2014, April 2, 2012 [cit. 2014-05-05]. 
	Dostupné z: \url{https://developers.google.com/protocol-buffers/}.

\bibitem{pbfenc}
	{\sc Google.}
	\emph{Encoding - Protocol Buffers} [online]. 
	2014, April 2, 2012 [cit. 2014-05-07]. 
	Dostupné z: \url{https://developers.google.com/protocol-buffers/docs/encoding}.

\bibitem{pbfspec}
	{\sc Google.}
	\emph{Language Guide - Protocol Buffers} [online]. 
	2014, April 22, 2014 [cit. 2014-05-07]. 
	Dostupné z:	\url{https://developers.google.com/protocol-buffers/docs/proto}.

\bibitem{lineint}
	{\sc Weisstein} Eric W.
	\emph{Line-Line Intersection.} [online]. 
	2014, May 12, 2014 [cit. 2014-05-12]. 
	Dostupné z:	\url{http://mathworld.wolfram.com/Line-LineIntersection.html}.

\bibitem{zametani}
	{\sc Mareš} Martin
	\emph{Geometrické algortimy} [online]. 
	2014, 17.11.2013 [cit. 2014-05-13]. 
	Dostupné z:	\url{http://mj.ucw.cz/vyuka/ads/43-geom.pdf}.
	

\end{thebibliography}


%%% Tabulky v bakalářské práci, existují-li.
\chapwithtoc{Seznam tabulek}

%%% Použité zkratky v bakalářské práci, existují-li, včetně jejich vysvětlení.
\chapwithtoc{Seznam použitých zkratek}

%%% Přílohy k bakalářské práci, existují-li (různé dodatky jako výpisy programů,
%%% diagramy apod.). Každá příloha musí být alespoň jednou odkazována z vlastního
%%% textu práce. Přílohy se číslují.
\chapwithtoc{Přílohy}

\openright
\end{document}
